\section{Full result figures }
\label{app:figures}
This appendix reproduces \emph{every} figure generated by the three configuration-driven
pipeline runs of the latest experiment configuration (\texttt{column}, \texttt{parkes}, \texttt{tvb}; run 2026-05-27)
over the six-mouse cohort. A representative subset appears in the main text; the complete set is
collected here for reference. Within each panel grid the two \textsc{Anderson} case mice precede the
four \textsc{Arnold} controls.

\subsection{Network-level \texorpdfstring{$\Dst$}{Dst} across the cohort}
\begin{figure}[H]\centering
\begin{subfigure}{0.32\textwidth}\centering
\includegraphics[width=\linewidth]{appendix/D_st_cohort_column.png}
\caption{\texttt{column}}
\end{subfigure}\hfill
\begin{subfigure}{0.32\textwidth}\centering
\includegraphics[width=\linewidth]{appendix/D_st_cohort_parkes.png}
\caption{\texttt{parkes}}
\end{subfigure}\hfill
\begin{subfigure}{0.32\textwidth}\centering
\includegraphics[width=\linewidth]{appendix/D_st_cohort_tvb.png}
\caption{\texttt{tvb}}
\end{subfigure}\hfill
\caption{Per-mouse network $\Dst=\mathrm{trace}(\Omega)/N$ for each normalisation scheme. Only \texttt{column} separates the strains; under \texttt{parkes}/\texttt{tvb} the cohort ranges overlap.}
\end{figure}

\subsection{Per-region stability centralities}
\begin{figure}[H]\centering
\begin{subfigure}{0.32\textwidth}\centering
\includegraphics[width=\linewidth]{appendix/stabilityCentralities_Anderson_1_column_figure.png}
\caption{\textsc{Anderson\_1}}
\end{subfigure}\hfill
\begin{subfigure}{0.32\textwidth}\centering
\includegraphics[width=\linewidth]{appendix/stabilityCentralities_Anderson_2_column_figure.png}
\caption{\textsc{Anderson\_2}}
\end{subfigure}\hfill
\begin{subfigure}{0.32\textwidth}\centering
\includegraphics[width=\linewidth]{appendix/stabilityCentralities_Arnold_2_column_figure.png}
\caption{\textsc{Arnold\_2}}
\end{subfigure}\hfill
\begin{subfigure}{0.32\textwidth}\centering
\includegraphics[width=\linewidth]{appendix/stabilityCentralities_Arnold_3_column_figure.png}
\caption{\textsc{Arnold\_3}}
\end{subfigure}\hfill
\begin{subfigure}{0.32\textwidth}\centering
\includegraphics[width=\linewidth]{appendix/stabilityCentralities_Arnold_4_column_figure.png}
\caption{\textsc{Arnold\_4}}
\end{subfigure}\hfill
\begin{subfigure}{0.32\textwidth}\centering
\includegraphics[width=\linewidth]{appendix/stabilityCentralities_Arnold_5_column_figure.png}
\caption{\textsc{Arnold\_5}}
\end{subfigure}\hfill
\caption{Per-region susceptibility $\Dto$ and influence $\Dfrom$ for all six mice, \texttt{column} scheme.}
\end{figure}
\begin{figure}[H]\centering
\includegraphics[width=0.8\textwidth]{appendix/stabilityCentralities_summary_overview_column.png}
\caption{Cohort overview of stability centralities, \texttt{column} scheme.}
\end{figure}

\begin{figure}[H]\centering
\begin{subfigure}{0.32\textwidth}\centering
\includegraphics[width=\linewidth]{appendix/stabilityCentralities_Anderson_1_parkes_figure.png}
\caption{\textsc{Anderson\_1}}
\end{subfigure}\hfill
\begin{subfigure}{0.32\textwidth}\centering
\includegraphics[width=\linewidth]{appendix/stabilityCentralities_Anderson_2_parkes_figure.png}
\caption{\textsc{Anderson\_2}}
\end{subfigure}\hfill
\begin{subfigure}{0.32\textwidth}\centering
\includegraphics[width=\linewidth]{appendix/stabilityCentralities_Arnold_2_parkes_figure.png}
\caption{\textsc{Arnold\_2}}
\end{subfigure}\hfill
\begin{subfigure}{0.32\textwidth}\centering
\includegraphics[width=\linewidth]{appendix/stabilityCentralities_Arnold_3_parkes_figure.png}
\caption{\textsc{Arnold\_3}}
\end{subfigure}\hfill
\begin{subfigure}{0.32\textwidth}\centering
\includegraphics[width=\linewidth]{appendix/stabilityCentralities_Arnold_4_parkes_figure.png}
\caption{\textsc{Arnold\_4}}
\end{subfigure}\hfill
\begin{subfigure}{0.32\textwidth}\centering
\includegraphics[width=\linewidth]{appendix/stabilityCentralities_Arnold_5_parkes_figure.png}
\caption{\textsc{Arnold\_5}}
\end{subfigure}\hfill
\caption{Per-region susceptibility $\Dto$ and influence $\Dfrom$ for all six mice, \texttt{parkes} scheme.}
\end{figure}
\begin{figure}[H]\centering
\includegraphics[width=0.8\textwidth]{appendix/stabilityCentralities_summary_overview_parkes.png}
\caption{Cohort overview of stability centralities, \texttt{parkes} scheme.}
\end{figure}

\begin{figure}[H]\centering
\begin{subfigure}{0.32\textwidth}\centering
\includegraphics[width=\linewidth]{appendix/stabilityCentralities_Anderson_1_tvb_figure.png}
\caption{\textsc{Anderson\_1}}
\end{subfigure}\hfill
\begin{subfigure}{0.32\textwidth}\centering
\includegraphics[width=\linewidth]{appendix/stabilityCentralities_Anderson_2_tvb_figure.png}
\caption{\textsc{Anderson\_2}}
\end{subfigure}\hfill
\begin{subfigure}{0.32\textwidth}\centering
\includegraphics[width=\linewidth]{appendix/stabilityCentralities_Arnold_2_tvb_figure.png}
\caption{\textsc{Arnold\_2}}
\end{subfigure}\hfill
\begin{subfigure}{0.32\textwidth}\centering
\includegraphics[width=\linewidth]{appendix/stabilityCentralities_Arnold_3_tvb_figure.png}
\caption{\textsc{Arnold\_3}}
\end{subfigure}\hfill
\begin{subfigure}{0.32\textwidth}\centering
\includegraphics[width=\linewidth]{appendix/stabilityCentralities_Arnold_4_tvb_figure.png}
\caption{\textsc{Arnold\_4}}
\end{subfigure}\hfill
\begin{subfigure}{0.32\textwidth}\centering
\includegraphics[width=\linewidth]{appendix/stabilityCentralities_Arnold_5_tvb_figure.png}
\caption{\textsc{Arnold\_5}}
\end{subfigure}\hfill
\caption{Per-region susceptibility $\Dto$ and influence $\Dfrom$ for all six mice, \texttt{tvb} scheme.}
\end{figure}
\begin{figure}[H]\centering
\includegraphics[width=0.8\textwidth]{appendix/stabilityCentralities_summary_overview_tvb.png}
\caption{Cohort overview of stability centralities, \texttt{tvb} scheme.}
\end{figure}

\subsection{Centrality correlations}
\begin{figure}[H]\centering
\begin{subfigure}{0.48\textwidth}\centering
\includegraphics[width=\linewidth]{appendix/centrality_corr_mean_column.png}
\caption{Cohort mean}
\end{subfigure}\hfill
\begin{subfigure}{0.48\textwidth}\centering
\includegraphics[width=\linewidth]{appendix/centrality_corr_bars_column.png}
\caption{$\Dto$, $\Dfrom$ vs classical}
\end{subfigure}
\caption{Centrality correlation summary, \texttt{column} scheme.}
\end{figure}
\begin{figure}[H]\centering
\begin{subfigure}{0.32\textwidth}\centering
\includegraphics[width=\linewidth]{appendix/centrality_corr_Anderson_1_column.png}
\caption{\textsc{Anderson\_1}}
\end{subfigure}\hfill
\begin{subfigure}{0.32\textwidth}\centering
\includegraphics[width=\linewidth]{appendix/centrality_corr_Anderson_2_column.png}
\caption{\textsc{Anderson\_2}}
\end{subfigure}\hfill
\begin{subfigure}{0.32\textwidth}\centering
\includegraphics[width=\linewidth]{appendix/centrality_corr_Arnold_2_column.png}
\caption{\textsc{Arnold\_2}}
\end{subfigure}\hfill
\begin{subfigure}{0.32\textwidth}\centering
\includegraphics[width=\linewidth]{appendix/centrality_corr_Arnold_3_column.png}
\caption{\textsc{Arnold\_3}}
\end{subfigure}\hfill
\begin{subfigure}{0.32\textwidth}\centering
\includegraphics[width=\linewidth]{appendix/centrality_corr_Arnold_4_column.png}
\caption{\textsc{Arnold\_4}}
\end{subfigure}\hfill
\begin{subfigure}{0.32\textwidth}\centering
\includegraphics[width=\linewidth]{appendix/centrality_corr_Arnold_5_column.png}
\caption{\textsc{Arnold\_5}}
\end{subfigure}\hfill
\caption{Per-mouse centrality correlation matrices, \texttt{column} scheme.}
\end{figure}

\begin{figure}[H]\centering
\begin{subfigure}{0.48\textwidth}\centering
\includegraphics[width=\linewidth]{appendix/centrality_corr_mean_parkes.png}
\caption{Cohort mean}
\end{subfigure}\hfill
\begin{subfigure}{0.48\textwidth}\centering
\includegraphics[width=\linewidth]{appendix/centrality_corr_bars_parkes.png}
\caption{$\Dto$, $\Dfrom$ vs classical}
\end{subfigure}
\caption{Centrality correlation summary, \texttt{parkes} scheme.}
\end{figure}
\begin{figure}[H]\centering
\begin{subfigure}{0.32\textwidth}\centering
\includegraphics[width=\linewidth]{appendix/centrality_corr_Anderson_1_parkes.png}
\caption{\textsc{Anderson\_1}}
\end{subfigure}\hfill
\begin{subfigure}{0.32\textwidth}\centering
\includegraphics[width=\linewidth]{appendix/centrality_corr_Anderson_2_parkes.png}
\caption{\textsc{Anderson\_2}}
\end{subfigure}\hfill
\begin{subfigure}{0.32\textwidth}\centering
\includegraphics[width=\linewidth]{appendix/centrality_corr_Arnold_2_parkes.png}
\caption{\textsc{Arnold\_2}}
\end{subfigure}\hfill
\begin{subfigure}{0.32\textwidth}\centering
\includegraphics[width=\linewidth]{appendix/centrality_corr_Arnold_3_parkes.png}
\caption{\textsc{Arnold\_3}}
\end{subfigure}\hfill
\begin{subfigure}{0.32\textwidth}\centering
\includegraphics[width=\linewidth]{appendix/centrality_corr_Arnold_4_parkes.png}
\caption{\textsc{Arnold\_4}}
\end{subfigure}\hfill
\begin{subfigure}{0.32\textwidth}\centering
\includegraphics[width=\linewidth]{appendix/centrality_corr_Arnold_5_parkes.png}
\caption{\textsc{Arnold\_5}}
\end{subfigure}\hfill
\caption{Per-mouse centrality correlation matrices, \texttt{parkes} scheme.}
\end{figure}

\begin{figure}[H]\centering
\begin{subfigure}{0.48\textwidth}\centering
\includegraphics[width=\linewidth]{appendix/centrality_corr_mean_tvb.png}
\caption{Cohort mean}
\end{subfigure}\hfill
\begin{subfigure}{0.48\textwidth}\centering
\includegraphics[width=\linewidth]{appendix/centrality_corr_bars_tvb.png}
\caption{$\Dto$, $\Dfrom$ vs classical}
\end{subfigure}
\caption{Centrality correlation summary, \texttt{tvb} scheme.}
\end{figure}
\begin{figure}[H]\centering
\begin{subfigure}{0.32\textwidth}\centering
\includegraphics[width=\linewidth]{appendix/centrality_corr_Anderson_1_tvb.png}
\caption{\textsc{Anderson\_1}}
\end{subfigure}\hfill
\begin{subfigure}{0.32\textwidth}\centering
\includegraphics[width=\linewidth]{appendix/centrality_corr_Anderson_2_tvb.png}
\caption{\textsc{Anderson\_2}}
\end{subfigure}\hfill
\begin{subfigure}{0.32\textwidth}\centering
\includegraphics[width=\linewidth]{appendix/centrality_corr_Arnold_2_tvb.png}
\caption{\textsc{Arnold\_2}}
\end{subfigure}\hfill
\begin{subfigure}{0.32\textwidth}\centering
\includegraphics[width=\linewidth]{appendix/centrality_corr_Arnold_3_tvb.png}
\caption{\textsc{Arnold\_3}}
\end{subfigure}\hfill
\begin{subfigure}{0.32\textwidth}\centering
\includegraphics[width=\linewidth]{appendix/centrality_corr_Arnold_4_tvb.png}
\caption{\textsc{Arnold\_4}}
\end{subfigure}\hfill
\begin{subfigure}{0.32\textwidth}\centering
\includegraphics[width=\linewidth]{appendix/centrality_corr_Arnold_5_tvb.png}
\caption{\textsc{Arnold\_5}}
\end{subfigure}\hfill
\caption{Per-mouse centrality correlation matrices, \texttt{tvb} scheme.}
\end{figure}

\subsection{Anderson vs Arnold per-region contrasts}
\begin{figure}[H]\centering
\begin{subfigure}{0.32\textwidth}\centering
\includegraphics[width=\linewidth]{appendix/compare_anderson_vs_arnold_Anderson_1_column_D_to_i.png}
\caption{\textsc{Anderson\_1}, $\Dto$}
\end{subfigure}\hfill
\begin{subfigure}{0.32\textwidth}\centering
\includegraphics[width=\linewidth]{appendix/compare_anderson_vs_arnold_Anderson_1_column_D_k_to.png}
\caption{\textsc{Anderson\_1}, $\Dfrom$}
\end{subfigure}\hfill
\begin{subfigure}{0.32\textwidth}\centering
\includegraphics[width=\linewidth]{appendix/compare_anderson_vs_arnold_Anderson_1_column_BC.png}
\caption{\textsc{Anderson\_1}, BC}
\end{subfigure}\hfill
\begin{subfigure}{0.32\textwidth}\centering
\includegraphics[width=\linewidth]{appendix/compare_anderson_vs_arnold_Anderson_2_column_D_to_i.png}
\caption{\textsc{Anderson\_2}, $\Dto$}
\end{subfigure}\hfill
\begin{subfigure}{0.32\textwidth}\centering
\includegraphics[width=\linewidth]{appendix/compare_anderson_vs_arnold_Anderson_2_column_D_k_to.png}
\caption{\textsc{Anderson\_2}, $\Dfrom$}
\end{subfigure}\hfill
\begin{subfigure}{0.32\textwidth}\centering
\includegraphics[width=\linewidth]{appendix/compare_anderson_vs_arnold_Anderson_2_column_BC.png}
\caption{\textsc{Anderson\_2}, BC}
\end{subfigure}\hfill
\caption{Each case mouse against the four-control band, \texttt{column} scheme: susceptibility $\Dto$, influence $\Dfrom$, and betweenness BC.}
\end{figure}

\begin{figure}[H]\centering
\begin{subfigure}{0.32\textwidth}\centering
\includegraphics[width=\linewidth]{appendix/compare_anderson_vs_arnold_Anderson_1_parkes_D_to_i.png}
\caption{\textsc{Anderson\_1}, $\Dto$}
\end{subfigure}\hfill
\begin{subfigure}{0.32\textwidth}\centering
\includegraphics[width=\linewidth]{appendix/compare_anderson_vs_arnold_Anderson_1_parkes_D_k_to.png}
\caption{\textsc{Anderson\_1}, $\Dfrom$}
\end{subfigure}\hfill
\begin{subfigure}{0.32\textwidth}\centering
\includegraphics[width=\linewidth]{appendix/compare_anderson_vs_arnold_Anderson_1_parkes_BC.png}
\caption{\textsc{Anderson\_1}, BC}
\end{subfigure}\hfill
\begin{subfigure}{0.32\textwidth}\centering
\includegraphics[width=\linewidth]{appendix/compare_anderson_vs_arnold_Anderson_2_parkes_D_to_i.png}
\caption{\textsc{Anderson\_2}, $\Dto$}
\end{subfigure}\hfill
\begin{subfigure}{0.32\textwidth}\centering
\includegraphics[width=\linewidth]{appendix/compare_anderson_vs_arnold_Anderson_2_parkes_D_k_to.png}
\caption{\textsc{Anderson\_2}, $\Dfrom$}
\end{subfigure}\hfill
\begin{subfigure}{0.32\textwidth}\centering
\includegraphics[width=\linewidth]{appendix/compare_anderson_vs_arnold_Anderson_2_parkes_BC.png}
\caption{\textsc{Anderson\_2}, BC}
\end{subfigure}\hfill
\caption{Each case mouse against the four-control band, \texttt{parkes} scheme: susceptibility $\Dto$, influence $\Dfrom$, and betweenness BC.}
\end{figure}

\begin{figure}[H]\centering
\begin{subfigure}{0.32\textwidth}\centering
\includegraphics[width=\linewidth]{appendix/compare_anderson_vs_arnold_Anderson_1_tvb_D_to_i.png}
\caption{\textsc{Anderson\_1}, $\Dto$}
\end{subfigure}\hfill
\begin{subfigure}{0.32\textwidth}\centering
\includegraphics[width=\linewidth]{appendix/compare_anderson_vs_arnold_Anderson_1_tvb_D_k_to.png}
\caption{\textsc{Anderson\_1}, $\Dfrom$}
\end{subfigure}\hfill
\begin{subfigure}{0.32\textwidth}\centering
\includegraphics[width=\linewidth]{appendix/compare_anderson_vs_arnold_Anderson_1_tvb_BC.png}
\caption{\textsc{Anderson\_1}, BC}
\end{subfigure}\hfill
\begin{subfigure}{0.32\textwidth}\centering
\includegraphics[width=\linewidth]{appendix/compare_anderson_vs_arnold_Anderson_2_tvb_D_to_i.png}
\caption{\textsc{Anderson\_2}, $\Dto$}
\end{subfigure}\hfill
\begin{subfigure}{0.32\textwidth}\centering
\includegraphics[width=\linewidth]{appendix/compare_anderson_vs_arnold_Anderson_2_tvb_D_k_to.png}
\caption{\textsc{Anderson\_2}, $\Dfrom$}
\end{subfigure}\hfill
\begin{subfigure}{0.32\textwidth}\centering
\includegraphics[width=\linewidth]{appendix/compare_anderson_vs_arnold_Anderson_2_tvb_BC.png}
\caption{\textsc{Anderson\_2}, BC}
\end{subfigure}\hfill
\caption{Each case mouse against the four-control band, \texttt{tvb} scheme: susceptibility $\Dto$, influence $\Dfrom$, and betweenness BC.}
\end{figure}

\subsection{Left--right asymmetry}
\begin{figure}[H]\centering
\begin{subfigure}{0.46\textwidth}\centering
\includegraphics[width=\linewidth]{appendix/compare_LR_asymmetry_Anderson_1_column.png}
\caption{\textsc{Anderson\_1}}
\end{subfigure}\hfill
\begin{subfigure}{0.46\textwidth}\centering
\includegraphics[width=\linewidth]{appendix/compare_LR_asymmetry_Anderson_2_column.png}
\caption{\textsc{Anderson\_2}}
\end{subfigure}\hfill
\begin{subfigure}{0.46\textwidth}\centering
\includegraphics[width=\linewidth]{appendix/LR_asymmetry_groupTest_column.png}
\caption{Homotopic-pair group test}
\end{subfigure}\hfill
\begin{subfigure}{0.46\textwidth}\centering
\includegraphics[width=\linewidth]{appendix/LR_asymmetry_systematic_column.png}
\caption{Systematic side bias}
\end{subfigure}
\caption{Left--right asymmetry views, \texttt{column} scheme.}
\end{figure}

\begin{figure}[H]\centering
\begin{subfigure}{0.46\textwidth}\centering
\includegraphics[width=\linewidth]{appendix/compare_LR_asymmetry_Anderson_1_parkes.png}
\caption{\textsc{Anderson\_1}}
\end{subfigure}\hfill
\begin{subfigure}{0.46\textwidth}\centering
\includegraphics[width=\linewidth]{appendix/compare_LR_asymmetry_Anderson_2_parkes.png}
\caption{\textsc{Anderson\_2}}
\end{subfigure}\hfill
\begin{subfigure}{0.46\textwidth}\centering
\includegraphics[width=\linewidth]{appendix/LR_asymmetry_groupTest_parkes.png}
\caption{Homotopic-pair group test}
\end{subfigure}\hfill
\begin{subfigure}{0.46\textwidth}\centering
\includegraphics[width=\linewidth]{appendix/LR_asymmetry_systematic_parkes.png}
\caption{Systematic side bias}
\end{subfigure}
\caption{Left--right asymmetry views, \texttt{parkes} scheme.}
\end{figure}

\begin{figure}[H]\centering
\begin{subfigure}{0.46\textwidth}\centering
\includegraphics[width=\linewidth]{appendix/compare_LR_asymmetry_Anderson_1_tvb.png}
\caption{\textsc{Anderson\_1}}
\end{subfigure}\hfill
\begin{subfigure}{0.46\textwidth}\centering
\includegraphics[width=\linewidth]{appendix/compare_LR_asymmetry_Anderson_2_tvb.png}
\caption{\textsc{Anderson\_2}}
\end{subfigure}\hfill
\begin{subfigure}{0.46\textwidth}\centering
\includegraphics[width=\linewidth]{appendix/LR_asymmetry_groupTest_tvb.png}
\caption{Homotopic-pair group test}
\end{subfigure}\hfill
\begin{subfigure}{0.46\textwidth}\centering
\includegraphics[width=\linewidth]{appendix/LR_asymmetry_systematic_tvb.png}
\caption{Systematic side bias}
\end{subfigure}
\caption{Left--right asymmetry views, \texttt{tvb} scheme.}
\end{figure}

\subsection{Cohort comparison figures (\texttt{column})}
These two summary figures also appear in the main text (Fig.~\ref{fig:outliers}, Fig.~\ref{fig:lr}).
\begin{figure}[H]\centering
\begin{subfigure}{0.49\textwidth}\centering
\includegraphics[width=\linewidth]{appendix/region_susc_outliers.png}
\caption{Region susceptibility outliers}
\end{subfigure}\hfill
\begin{subfigure}{0.49\textwidth}\centering
\includegraphics[width=\linewidth]{appendix/laterality_influence_outliers.png}
\caption{Laterality influence outliers}
\end{subfigure}
\caption{Cohort-level Bonferroni-significant outliers, \texttt{column} scheme.}
\end{figure}

\section{Full result tables}
\label{app:tables}
Numeric summaries underlying the figures above. Complete per-node CSVs for every run are written
to the run folders by the pipeline (see Appendix~\ref{app:repro}).

\subsection{Network \texorpdfstring{$\Dst$}{Dst} and spectral radius \texorpdfstring{$\rhoC$}{rho(C)}}
\begin{table}[H]\centering\small
\caption{Per-mouse $\Dst$ and $\rhoC$ for the three schemes. Case mice (\textsc{Anderson}) above the rule.}
\begin{tabular}{@{}l ccc ccc@{}}\toprule
& \multicolumn{3}{c}{$\Dst$} & \multicolumn{3}{c}{$\rhoC$} \\
\cmidrule(lr){2-4}\cmidrule(lr){5-7}
Mouse & \texttt{column} & \texttt{parkes} & \texttt{tvb} & \texttt{column} & \texttt{parkes} & \texttt{tvb} \\\midrule
Anderson\_1 & 0.7924 & 653.3 & 271.9 & 0.95 & 0.999988 & 0.999152 \\
Anderson\_2 & 0.8069 & 545.2 & 274.9 & 0.95 & 0.999986 & 0.999152 \\
\midrule
Arnold\_2 & 0.7582 & 528.9 & 272.1 & 0.95 & 0.999986 & 0.999152 \\
Arnold\_3 & 0.7714 & 643.9 & 272.1 & 0.95 & 0.999988 & 0.999152 \\
Arnold\_4 & 0.7721 & 686.9 & 268.7 & 0.95 & 0.999989 & 0.999152 \\
Arnold\_5 & 0.7884 & 635.8 & 269.5 & 0.95 & 0.999988 & 0.999152 \\
\bottomrule\end{tabular}\end{table}

\subsection{Top regions by stability susceptibility}
\begin{table}[H]\centering\small
\caption{Top 15 non-trivial regions by mean susceptibility rank, \texttt{column} scheme (lower rank = more central).}
\begin{tabular}{@{}rlcc@{}}\toprule
\# & Region & mean rank $\Dto$ & mean rank $\Dfrom$ \\\midrule
1 & L-CORTEX\_SOMATOSENSORY & 3 & 1.67 \\
2 & R-CORTEX\_SOMATOSENSORY & 4.33 & 3.83 \\
3 & L-CEREBELLUM & 5.5 & 8.5 \\
4 & R-CEREBELLUM & 8.17 & 11.2 \\
5 & R-MEDULLA & 9 & 16 \\
6 & L-CORTEX\_VISUAL & 11.3 & 6.83 \\
7 & L-CORTEX\_MOTOR & 12 & 14.7 \\
8 & R-THALAMUS & 12 & 10.8 \\
9 & L-MEDULLA & 14.3 & 20.3 \\
10 & L-OLFACTORY & 14.5 & 9.33 \\
11 & L-PARAHIPPOCAMPAL/ENTORHINAL & 15 & 16.5 \\
12 & R-CORTEX\_VISUAL & 15.5 & 9.67 \\
13 & L-BASALGANGLIA\_STRIATUM & 16.2 & 12.7 \\
14 & R-OLFACTORY & 16.5 & 7.33 \\
15 & R-CORTEX\_MOTOR & 16.7 & 13.3 \\
\bottomrule\end{tabular}\end{table}
\begin{table}[H]\centering\small
\caption{Top 15 non-trivial regions by mean susceptibility rank, \texttt{parkes} scheme (lower rank = more central).}
\begin{tabular}{@{}rlcc@{}}\toprule
\# & Region & mean rank $\Dto$ & mean rank $\Dfrom$ \\\midrule
1 & L-CORTEX\_SOMATOSENSORY & 2.17 & 2.17 \\
2 & R-CORTEX\_SOMATOSENSORY & 4.67 & 4.67 \\
3 & L-CORTEX\_VISUAL & 8.33 & 8.33 \\
4 & L-OLFACTORY & 8.67 & 8.67 \\
5 & R-OLFACTORY & 8.83 & 8.83 \\
6 & L-WHITE\_MATTER & 9.17 & 9.17 \\
7 & L-BASALGANGLIA\_STRIATUM & 9.67 & 9.67 \\
8 & R-WHITE\_MATTER & 10.5 & 10.5 \\
9 & L-CORTEX\_MOTOR & 11.3 & 11.3 \\
10 & R-CORTEX\_VISUAL & 11.5 & 11.5 \\
11 & L-CEREBELLUM & 13 & 13 \\
12 & R-THALAMUS & 13 & 13 \\
13 & R-BASALGANGLIA\_STRIATUM & 13.2 & 13.2 \\
14 & R-CEREBELLUM & 13.8 & 13.8 \\
15 & L-PARAHIPPOCAMPAL/ENTORHINAL & 14 & 14 \\
\bottomrule\end{tabular}\end{table}
\begin{table}[H]\centering\small
\caption{Top 15 non-trivial regions by mean susceptibility rank, \texttt{tvb} scheme (lower rank = more central).}
\begin{tabular}{@{}rlccc@{}}\toprule
\# & Region & mean rank $\Dto$ & mean rank $\Dfrom$ & mean rank $\xcrit$ \\\midrule
1 & R-THALAMUS\_AUDITORY & 5.17 & 5.17 & 1.83 \\
2 & L-THALAMUS\_AUDITORY & 5.67 & 5.67 & 2.17 \\
3 & L-BASALGANGLIA\_PALLIDUM & 9 & 9 & 4.83 \\
4 & R-THALAMUS\_VISUAL & 9.67 & 9.67 & 5.67 \\
5 & R-BASALGANGLIA\_PALLIDUM & 10 & 10 & 6.17 \\
6 & L-THALAMUS\_VISUAL & 10.7 & 10.7 & 6.33 \\
7 & R-EPITHALAMUS/HABENULA & 11.8 & 11.8 & 7.5 \\
8 & R-LIMBIC\_CORTEX & 12.2 & 12.2 & 7.83 \\
9 & L-LIMBIC\_CORTEX & 12.8 & 12.8 & 8.33 \\
10 & L-EPITHALAMUS/HABENULA & 13.7 & 13.7 & 9.17 \\
11 & L-CORTEX\_GUSTATORY & 15.8 & 15.8 & 11.2 \\
12 & R-CORTEX\_GUSTATORY & 17 & 17 & 12.2 \\
13 & R-SEPTUM & 17.5 & 17.5 & 12.5 \\
14 & L-SEPTUM & 19.3 & 19.3 & 14.2 \\
15 & L-BASAL\_FOREBRAIN & 20.8 & 20.8 & 15 \\
\bottomrule\end{tabular}\end{table}

\subsection{Anderson vs Arnold per-region outliers}
\begin{table}[H]\centering\small
\caption{Flagged per-region contrasts ($|z|\ge2$) vs the control band, \texttt{column} scheme. $p_b$ = Bonferroni-corrected $p$; $^{\ast}$ marks significance at $\alpha=0.05$.}
\begin{tabular}{@{}llrrrrrr@{}}\toprule
Mouse & Region & $z_{\mathrm{susc}}$ & $p_b^{\mathrm{susc}}$ & $z_{\mathrm{infl}}$ & $p_b^{\mathrm{infl}}$ & $z_{\mathrm{BC}}$ & $p_b^{\mathrm{BC}}$ \\\midrule
Anderson\_1 & R-CRANIAL\_NERVES & 1.12 & 1 & -0.24 & 1 & 119.50 & $<\!10^{-15}$\textsuperscript{$\ast$} \\
Anderson\_1 & L-CRANIAL\_NERVES & 0.51 & 1 & -0.17 & 1 & 7.17 & $1.3\!\times\!10^{-10}$\textsuperscript{$\ast$} \\
Anderson\_1 & R-THALAMUS & -1.04 & 1 & -0.70 & 1 & 5.65 & $2.8\!\times\!10^{-6}$\textsuperscript{$\ast$} \\
Anderson\_1 & R-HIPPOCAMPAL FORMATION & 5.38 & $1.3\!\times\!10^{-5}$\textsuperscript{$\ast$} & 0.90 & 1 & -0.66 & 1 \\
Anderson\_1 & L-THALAMUS\_VISUAL & 0.35 & 1 & 4.42 & 0.00172\textsuperscript{$\ast$} & 1.50 & 1 \\
Anderson\_1 & R-BASALGANGLIA\_STRIATUM & 4.25 & 0.00379\textsuperscript{$\ast$} & -0.90 & 1 & -0.69 & 1 \\
Anderson\_1 & L-CORTEX\_GUSTATORY & 4.00 & 0.0111\textsuperscript{$\ast$} & -0.32 & 1 & -0.50 & 1 \\
Anderson\_2 & L-CRANIAL\_NERVES & 1.88 & 1 & -0.86 & 1 & 29.42 & $<\!10^{-15}$\textsuperscript{$\ast$} \\
Anderson\_2 & R-VENTRICLES\_CSF & 1.33 & 1 & -2.32 & 1 & 10.77 & $<\!10^{-15}$\textsuperscript{$\ast$} \\
Anderson\_2 & L-CORTEX\_GUSTATORY & 7.87 & $6.3\!\times\!10^{-13}$\textsuperscript{$\ast$} & 0.12 & 1 & -0.50 & 1 \\
Anderson\_2 & L-HIPPOCAMPAL FORMATION & 6.75 & $2.7\!\times\!10^{-9}$\textsuperscript{$\ast$} & 0.53 & 1 & 1.29 & 1 \\
Anderson\_2 & R-THALAMUS & -0.15 & 1 & -0.12 & 1 & 4.80 & $2.9\!\times\!10^{-4}$\textsuperscript{$\ast$} \\
Anderson\_2 & L-CORTEX\_TEMPORAL\_ASSOC & 4.57 & $8.8\!\times\!10^{-4}$\textsuperscript{$\ast$} & -0.91 & 1 & -0.70 & 1 \\
Anderson\_2 & L-CORTEX\_SOMATOSENSORY & 3.74 & 0.033\textsuperscript{$\ast$} & 1.31 & 1 & 0.07 & 1 \\
Anderson\_2 & L-WHITE\_MATTER & 3.03 & 0.433 & 0.77 & 1 & 3.65 & 0.0469\textsuperscript{$\ast$} \\
\bottomrule\end{tabular}\end{table}
\begin{table}[H]\centering\small
\caption{Flagged per-region contrasts ($|z|\ge2$) vs the control band, \texttt{parkes} scheme. $p_b$ = Bonferroni-corrected $p$; $^{\ast}$ marks significance at $\alpha=0.05$.}
\begin{tabular}{@{}llrrrrrr@{}}\toprule
Mouse & Region & $z_{\mathrm{susc}}$ & $p_b^{\mathrm{susc}}$ & $z_{\mathrm{infl}}$ & $p_b^{\mathrm{infl}}$ & $z_{\mathrm{BC}}$ & $p_b^{\mathrm{BC}}$ \\\midrule
Anderson\_1 & R-AMYGDALA & 7.40 & $2.2\!\times\!10^{-11}$\textsuperscript{$\ast$} & 7.40 & $2.2\!\times\!10^{-11}$\textsuperscript{$\ast$} & -- & 1 \\
Anderson\_1 & L-THALAMUS\_VISUAL & 3.86 & 0.0174\textsuperscript{$\ast$} & 3.86 & 0.0174\textsuperscript{$\ast$} & -- & 1 \\
Anderson\_2 & R-THALAMUS & -7.45 & $1.5\!\times\!10^{-11}$\textsuperscript{$\ast$} & -7.45 & $1.5\!\times\!10^{-11}$\textsuperscript{$\ast$} & 2.10 & 1 \\
Anderson\_2 & L-CORTEX\_MOTOR & 3.51 & 0.0692 & 3.51 & 0.0692 & 7.43 & $1.7\!\times\!10^{-11}$\textsuperscript{$\ast$} \\
\bottomrule\end{tabular}\end{table}
\begin{table}[H]\centering\small
\caption{Flagged per-region contrasts ($|z|\ge2$) vs the control band, \texttt{tvb} scheme. $p_b$ = Bonferroni-corrected $p$; $^{\ast}$ marks significance at $\alpha=0.05$.}
\begin{tabular}{@{}llrrrrrr@{}}\toprule
Mouse & Region & $z_{\mathrm{susc}}$ & $p_b^{\mathrm{susc}}$ & $z_{\mathrm{infl}}$ & $p_b^{\mathrm{infl}}$ & $z_{\mathrm{BC}}$ & $p_b^{\mathrm{BC}}$ \\\midrule
Anderson\_1 & R-MEDULLA & -0.67 & 1 & -0.67 & 1 & 6.78 & $1.8\!\times\!10^{-9}$\textsuperscript{$\ast$} \\
Anderson\_1 & R-WHITE\_MATTER & 4.15 & 0.00527\textsuperscript{$\ast$} & 4.15 & 0.00527\textsuperscript{$\ast$} & -1.01 & 1 \\
Anderson\_1 & L-THALAMUS\_VISUAL & -3.98 & 0.011\textsuperscript{$\ast$} & -3.98 & 0.011\textsuperscript{$\ast$} & -- & 1 \\
Anderson\_2 & R-CORTEX\_GUSTATORY & 7.29 & $4.9\!\times\!10^{-11}$\textsuperscript{$\ast$} & 7.29 & $4.9\!\times\!10^{-11}$\textsuperscript{$\ast$} & -- & 1 \\
Anderson\_2 & R-CORTEX\_MOTOR & -0.79 & 1 & -0.79 & 1 & 3.70 & 0.0345\textsuperscript{$\ast$} \\
\bottomrule\end{tabular}\end{table}

\subsection{Systematic laterality}
\begin{table}[H]\centering\small
\caption{Cohort-wide sign-rank test of the normalised laterality index $\mathrm{LI}_{\mathrm{norm}}$ against zero, all schemes. Side = inferred dominant hemisphere (nominal labels).}
\begin{tabular}{@{}lllrrrl@{}}\toprule
Scheme & Cohort & Metric & mean LI & sd LI & sign-rank $p$ & Side \\\midrule
\texttt{column} & Anderson & D\_susc & 0.00879 & 0.0188 & $7.9\!\times\!10^{-4}$ & L \\
\texttt{column} & Arnold & D\_susc & 0.00173 & 0.0131 & 0.0724 & none \\
\texttt{column} & Anderson & D\_infl & 0.0399 & 0.0982 & 0.0026 & L \\
\texttt{column} & Arnold & D\_infl & 0.00862 & 0.083 & 0.178 & none \\
\texttt{column} & Anderson & BC & 0.156 & 0.694 & 0.189 & none \\
\texttt{column} & Arnold & BC & -0.0734 & 0.668 & 0.291 & none \\
\midrule
\texttt{parkes} & Anderson & D\_susc & 0.256 & 0.435 & $1.0\!\times\!10^{-4}$ & L \\
\texttt{parkes} & Arnold & D\_susc & 0.0654 & 0.359 & 0.0203 & L \\
\texttt{parkes} & Anderson & D\_infl & 0.256 & 0.435 & $1.0\!\times\!10^{-4}$ & L \\
\texttt{parkes} & Arnold & D\_infl & 0.0654 & 0.359 & 0.0203 & L \\
\texttt{parkes} & Anderson & BC & 0.387 & 0.511 & 0.00609 & L \\
\texttt{parkes} & Arnold & BC & -0.0651 & 0.744 & 0.417 & none \\
\midrule
\texttt{tvb} & Anderson & D\_susc & -0.0332 & 0.0769 & 0.00151 & R \\
\texttt{tvb} & Arnold & D\_susc & -0.00186 & 0.0838 & 0.613 & none \\
\texttt{tvb} & Anderson & D\_infl & -0.0332 & 0.0769 & 0.00151 & R \\
\texttt{tvb} & Arnold & D\_infl & -0.00186 & 0.0838 & 0.613 & none \\
\texttt{tvb} & Anderson & BC & 0.292 & 0.622 & 0.033 & L \\
\texttt{tvb} & Arnold & BC & -0.0203 & 0.676 & 0.851 & none \\
\bottomrule\end{tabular}\end{table}

\subsection{Homotopic-pair laterality tests}
\begin{table}[H]\centering\small
\caption{Per-pair Welch $t$-test between strains, \texttt{column} scheme: 10 smallest Bonferroni $p$. $\checkmark$ marks $p_b<0.05$.}
\begin{tabular}{@{}llrrrrc@{}}\toprule
Region & Metric & $\mathrm{LI}_{\mathrm{And}}$ & $\mathrm{LI}_{\mathrm{Arn}}$ & $t$ & $p_b$ & sig. \\\midrule
CORTEX\_VISUAL & D\_susc & 0.016 & 0.00737 & 0.58 & 1 &  \\
CORTEX\_SOMATOSENSORY & D\_susc & 0.0119 & 0.0146 & -0.10 & 1 &  \\
CORTEX\_MOTOR & D\_susc & 0.0261 & -0.000446 & 2.45 & 1 &  \\
CORTEX\_AUDITORY & D\_susc & 0.0213 & 0.0125 & 0.36 & 1 &  \\
CORTEX\_TEMPORAL\_ASSOC & D\_susc & 0.0252 & 0.00942 & 3.27 & 1 &  \\
CORTEX\_PREFRONTAL & D\_susc & 0.00152 & 0.00215 & -0.11 & 1 &  \\
CORTEX\_INSULAR & D\_susc & 0.022 & -0.00112 & 1.01 & 1 &  \\
CORTEX\_GUSTATORY & D\_susc & 0.0486 & -0.00299 & 2.68 & 1 &  \\
CORTEX\_RETROSPLENIAL & D\_susc & 0.00464 & -0.00525 & 1.12 & 1 &  \\
PARAHIPPOCAMPAL/ENTORHINAL & D\_susc & 0.0274 & -0.000259 & 1.46 & 1 &  \\
\bottomrule\end{tabular}\end{table}
\begin{table}[H]\centering\small
\caption{Per-pair Welch $t$-test between strains, \texttt{parkes} scheme: 10 smallest Bonferroni $p$. $\checkmark$ marks $p_b<0.05$.}
\begin{tabular}{@{}llrrrrc@{}}\toprule
Region & Metric & $\mathrm{LI}_{\mathrm{And}}$ & $\mathrm{LI}_{\mathrm{Arn}}$ & $t$ & $p_b$ & sig. \\\midrule
CORTEX\_VISUAL & D\_susc & 0.533 & 0.0803 & 0.98 & 1 &  \\
CORTEX\_SOMATOSENSORY & D\_susc & 0.17 & 0.223 & -0.08 & 1 &  \\
CORTEX\_MOTOR & D\_susc & 0.256 & -0.0398 & 0.51 & 1 &  \\
CORTEX\_AUDITORY & D\_susc & -0.00396 & 0.328 & -0.50 & 1 &  \\
CORTEX\_TEMPORAL\_ASSOC & D\_susc & 0.732 & 0.219 & 1.99 & 1 &  \\
CORTEX\_PREFRONTAL & D\_susc & 0.256 & 0.117 & 1.18 & 1 &  \\
CORTEX\_INSULAR & D\_susc & 0.321 & -0.0774 & 0.57 & 1 &  \\
CORTEX\_GUSTATORY & D\_susc & 0.298 & -0.107 & 0.58 & 1 &  \\
CORTEX\_RETROSPLENIAL & D\_susc & 0.523 & 0.0136 & 1.48 & 1 &  \\
PARAHIPPOCAMPAL/ENTORHINAL & D\_susc & 0.887 & 0.337 & 2.41 & 1 &  \\
\bottomrule\end{tabular}\end{table}
\begin{table}[H]\centering\small
\caption{Per-pair Welch $t$-test between strains, \texttt{tvb} scheme: 10 smallest Bonferroni $p$. $\checkmark$ marks $p_b<0.05$.}
\begin{tabular}{@{}llrrrrc@{}}\toprule
Region & Metric & $\mathrm{LI}_{\mathrm{And}}$ & $\mathrm{LI}_{\mathrm{Arn}}$ & $t$ & $p_b$ & sig. \\\midrule
CORTEX\_VISUAL & D\_susc & -0.0672 & 0.00111 & -0.67 & 1 &  \\
CORTEX\_SOMATOSENSORY & D\_susc & 0.0358 & -0.0212 & 1.45 & 1 &  \\
CORTEX\_MOTOR & D\_susc & 0.021 & 0.0605 & -0.61 & 1 &  \\
CORTEX\_AUDITORY & D\_susc & 0.0934 & -0.0349 & 1.88 & 1 &  \\
CORTEX\_TEMPORAL\_ASSOC & D\_susc & -0.116 & -0.000611 & -4.30 & 1 &  \\
CORTEX\_PREFRONTAL & D\_susc & -0.0413 & -0.011 & -1.61 & 1 &  \\
CORTEX\_INSULAR & D\_susc & -0.0333 & 0.0562 & -1.09 & 1 &  \\
CORTEX\_GUSTATORY & D\_susc & -0.0211 & 0.0412 & -0.76 & 1 &  \\
CORTEX\_RETROSPLENIAL & D\_susc & -0.149 & -0.0104 & -1.04 & 1 &  \\
PARAHIPPOCAMPAL/ENTORHINAL & D\_susc & -0.158 & -0.0639 & -1.12 & 1 &  \\
\bottomrule\end{tabular}\end{table}

